%\documentclass[12pt,oneside]{amsart}
\documentclass{scrartcl}

%\documentclass{article}

\usepackage{amsthm,amsmath,amssymb}
\usepackage[numbers]{natbib}

%\usepackage[colorlinks]{hyperref}
\usepackage{hyperref}
\hypersetup{
    colorlinks,%
    citecolor=black,%
    filecolor=black,%
    linkcolor=black,%
    urlcolor=black
}
\usepackage{hypernat}
\usepackage[top=2.5cm, bottom=2.5cm, left=2.8cm,
right=2.8cm]{geometry}
\usepackage{graphicx}


%\setlength{\parskip}{0pt}
%\setlength{\parsep}{0pt}
%\setlength{\headsep}{0pt}
%\setlength{\topskip}{0pt}
%\setlength{\topmargin}{0pt}
%\setlength{\topsep}{0pt}
%\setlength{\partopsep}{0pt}

\linespread{1}

%\usepackage{marvosym}

%\textwidth = 6.5 in \textheight = 9.2 in \oddsidemargin = 0.0 in
%\evensidemargin = 0.0 in \topmargin = -0.0 in \headheight = 0.0 in
%\headsep = 0.0in \parskip = 0.0in \parindent = 0.0in
%\def\baselinestretch{0.871}


\newtheorem{theorem}{Theorem}[section]
\newtheorem{proposition}[theorem]{Proposition}
\newtheorem{problem}[theorem]{Problem}
\newtheorem{fact}[theorem]{Fact}
\newtheorem{lemma}[theorem]{Lemma}
\newtheorem{defn}[theorem]{Definition}
\newtheorem{corollary}[theorem]{Corollary}
\newtheorem{remark}{Remark}[section]
\newtheorem{example}[theorem]{Example}

\newcommand{\E}{{\mathbb E}}
\newcommand{\se}{{\mathcal{E}}}
\newcommand{\R}{{\mathbb R}}
\renewcommand{\P}{{\mathbb P}}
\newcommand{\ac}{{\mathcal{A}}}
\newcommand{\A}{{\mathcal{A}}}
\newcommand{\B}{{\mathcal{B}}}
\newcommand{\C}{{\mathcal{C}}}
\newcommand{\M}{{\mathcal{M}}}
\newcommand{\Mf}{{\mathfrak{M}}}
\newcommand{\T}{{\mathcal{T}}}
\newcommand{\Z}{{\mathcal{Z}}}
\newcommand{\K}{{\mathcal{K}}}
\newcommand{\lc}{{\mathcal{L}}}
\newcommand{\Ps}{{\mathcal{P}}}
\newcommand{\G}{{\mathcal{G}}}
\newcommand{\pa}{{\mathring{p}}}
\newcommand{\F}{{\cal F}}
\newcommand{\samp}{{\mathcal{X}}}
\newcommand{\X}{{\mathcal{X}}}
\newcommand{\hi}{{\hat{\Phi}}}
\newcommand{\data}{{\text{data}}}
\newcommand{\relu}{{\text{ReLU}}}

\newcommand{\ridge}{{\hat{\beta}_{\text{ridge}}(\lambda)}}
\newcommand{\lasso}{{\hat{\beta}_{\text{lasso}}(\lambda)}}
\newcommand{\ridgemu}{{\hat{\mu}_t^{\text{ridge}}(\lambda)}}
\newcommand{\lassomu}{{\hat{\mu}_t^{\text{lasso}}(\lambda)}}


\newcounter{rcnt}[section]
\renewcommand{\thercnt}{(\roman{rcnt})}

\def\qt#1{\qquad\text{#1}}

\def\argmin{\mathop{\rm argmin}}
\def\argmax{\mathop{\rm argmax}}


\setlength{\parskip}{1.4 \medskipamount}
%\sloppy
%\linespread{1.3}

\begin{document}

\title{STAT 238 - Bayesian Statistics  \\
  Lecture Twenty One} 
\subtitle{Spring 2026, UC Berkeley}

\author{Aditya Guntuboyina}


\date{13 March 2026}

\maketitle

\section{Recap: Model from past few lectures}

We studied the following model in the past few lectures.

We have a response variable $y$ and a single covariate $x$. In our
example, $y$ denotes weekly earnings and $x$ denotes years of
experience. The covariate $x$ takes the values $0, 1, \dots, m$ for
some fixed integer $m$. 

Our data is $(x_i, y_i), i = 1, \dots, n$. The model is:
\begin{align}\label{disc_mod}
  y_i = \beta_0 + \beta_1 x_i + \beta_2 \relu(x_i - 1) + \dots +
  \beta_{m}\relu(x_i - (m-1)) + \epsilon_i
\end{align}
where $\relu(u) := \max(u, 0)$, and $\epsilon_i
\overset{\text{i.i.d}}{\sim} N(0, \sigma^2)$. Note we did not include
$\relu(x_i - m)$ because it always equals 0.

We discussed Bayesian inference in this model using the prior: 
\begin{align}\label{coef_prior}
  \beta_0 \sim N(0, C), \beta_1 \sim N(0, C), \beta_2, \dots, \beta_m
  \overset{\text{i.i.d}}{\sim} N(0, \tau^2)
\end{align}
Towards the end of last lecture, we remarked that the resulting model
on the regression function: 
\begin{align*}
  f(x) = \beta_0 + \beta_1 x + \beta_2 \relu(x - 1) + \dots +
  \beta_{m}\relu(x - (m-1))
\end{align*}
with prior \eqref{coef_prior} is very similar to:
\begin{align}\label{ibm_mod}
  f(x) = \beta_0 + \beta_1 x + \tau I(x)
\end{align}
where $I(x)$ is integrated Brownian motion (and $\beta_0, \beta_1
\overset{\text{i.i.d}}{\sim} N(0, C)$). We shall some intuition today
as to why the Integrated Brownian Motion is appearing here. 

\section{Brownian Motion and Integrated Brownian Motion}
Before defining Brownian motion, let us recall the definition of
Brownian motion. Brownian motion on $[0, M]$ (for some fixed $M > 0$)
is a stochastic process $\{B(t), 0 \le t \le M\}$ characterized by the
following two properties:

\begin{enumerate}
\item Every realization is a continuous function on $[0,M]$, and
  $B(0) = 0$.

\item For every fixed $t_1,\dots,t_k \in [0,M]$, the random vector
  $(B(t_1),\dots,B(t_k))$ has a multivariate normal distribution with
  mean zero and covariance matrix
  \[
  \text{Cov}(B(t_i), B(t_j)) = \min(t_i,t_j).
  \]
\end{enumerate}

This definition suggests the following method for simulating
$B(t)$. Construct the dense grid
\[
0 = t_0 < t_1 < \dots < t_N = M
\]
on $[0,M]$ with $t_i = iM/N$. One can then simulate
$(B(t_1),\dots,B(t_N))$ from the multivariate normal distribution with
mean zero and covariance matrix $\min(t_i,t_j)$.

However, this approach is computationally inefficient for large $N$.
Sampling directly from an $N \times N$ covariance matrix typically
requires a matrix factorization (such as a Cholesky decomposition),
which has computational cost on the order of $O(N^3)$.

For simulations in particular, the following alternative definition is
more appealing:
\begin{enumerate}
\item Every realization is a continuous function on $[0, M]$ with
  $B(0) = 0$.
\item The process has independent increments which means that $B(t_1),
  B(t_2) - B(t_1), \dots, B(t_k) - B(t_{k-1})$ are independent
  whenever $0 \leq t_1 < \dots < t_k \leq M$.
\item $B(t) - B(s) \sim N(0, t-s)$ whenever $0 \leq s \leq t$. 
\end{enumerate}
With this definition, we can simulate $B(t_1), \dots, B(t_N)$ for $t_i
= iM/N$ in the following way:
\begin{itemize}
\item First generate $\alpha_1, \dots, \alpha_N
  \overset{\text{i.i.d}}{\sim} N(0, M/N)$
\item Take $B(t_i) = \alpha_1 + \dots + \alpha_i$ for $i = 1, \dots,
  N$. 
\end{itemize}
Note that in this way $\alpha_i = B(t_i) - B(t_{i-1})$. The equation
$B(t_i) = \alpha_1 + \dots + \alpha_i$ can also be written as
\begin{align*}
  B(t) = \sum_{i=1}^N \alpha_i I\{t \geq t_i\} ~~ \text{ for $t \in
  \{t_1, \dots, t_N\}$}. 
\end{align*}
Because of the denseness of the grid (when $N$ is  large), we are
claiming that
\begin{align}\label{app_bm}
  B(t) \approx \sum_{i=1}^N \alpha_i I\{t \geq t_i\} ~~ \text{ for all
  $t \in [0, M]$}. 
\end{align}
This means that Brownian motion can be well-approximated (when $N$ is
large) by a piecewise constant process which jumps at each grid point
$t_i = iM/N$ by a small amount given by $N(0, M/N)$. 

We are now ready to define the Integrated Brownian Motion. This is
given by:
\begin{align}\label{ibm_def}
  I(t) = \int_0^t B(s) ds. 
\end{align}
Using the approximation \eqref{app_bm}, we can write
\begin{align*}
  I(t) \approx \int_0^t \left(\sum_{i=1}^N \alpha_i I\{t \geq t_i\} \right)
  ds = \sum_{i=1}^N \alpha_i \int_0^t I\{s \geq t_i\} ds =
  \sum_{i=1}^N \alpha_i (t - t_i)_+. 
\end{align*}
Therefore Integrated Brownian Motion can be approximated as:
\begin{align}\label{ibm_approx}
  I(t) \approx   \sum_{i=1}^N \alpha_i (t - t_i)_+ 
\end{align}
when $N$ is large, $t_i = iM/N$ and $\alpha_i
\overset{\text{i.i.d}}{\sim} N(0, M/N)$.

\section{The Model \eqref{ibm_mod}}
Using \eqref{ibm_approx}, the (prior) model \eqref{ibm_mod} can
therefore be written as:
\begin{align*}
  f(x) \approx \beta_0 + \beta_1 x + \tau \sum_{i=1}^N \alpha_i
  \left(x - \frac{iM}{N}\right)_+  
\end{align*}
Here $N$ is a large number. If we do a further coarse approximation
with $N = M$, we get back the model \eqref{disc_mod} if we make the
identification $\beta_i = \tau \alpha_{i-1}$ for $i = 2, \dots, N$ and $M =
m$. Therefore, the model \eqref{disc_mod} can be understood as a
coarse discretization of the model \eqref{ibm_mod} which stipulates
that $f$ is an Integrated Brownian Motion (scaled by $\tau$) and added
to a linear function $\beta_0 + \beta_1 x$ with the uniformative $N(0,
C)$ prior on the coefficients $\beta_0, \beta_1$.

The prior \eqref{ibm_mod} is an example of a Gaussian process prior. 

\section{Gaussian Process Regression}
Consider the usual nonparametric regression problem where the goal is
to estimate an unknown function $f : \Omega \rightarrow \R$ from
observations $(x_1, y_1), \dots, (x_n, y_n)$ under the model:
\begin{equation*}
  y_i = f(x_i) + \epsilon_i ~~\text{where $\epsilon_i
    \overset{\text{i.i.d}}{\sim} N(0, \sigma^2)$}. 
\end{equation*}
Here $\Omega$ denotes the domain which is a subset of $\R^d$ for some
$d \geq 1$. 

We assume that $\{f(x), x \in
\Omega\}$ forms a Gaussian process with mean function $\mu(x), x \in
\Omega$ and covariance 
function or \textit{kernel} given by $K(x, x')$ i.e.,
\begin{equation*}
  \text{Cov}(f(x), f(x')) = K(x, x') ~~\text{for all $x, x' \in
    \Omega$}. 
\end{equation*}
The kernel is positive semi-definite i.e., for every $N \geq 1$,
distinct points $u_1, \dots, u_N \in \Omega$, the $N \times N$
matrix with $(i, j)^{th}$ entry $K(u_i, u_j)$ is positive
semi-definite. Often, this matrix will be positive definite, and hence
invertible.

The mean function is usually taken to be zero. Here are some standard
special cases.
Here are some examples of Gaussian processes and kernels:
\begin{enumerate}
\item \textbf{Brownian Motion}: Here $\Omega = [0, \infty)$ and $K(s,
  t) = \min(s, t)$.
\item \textbf{(scaled) Brownian Motion plus constant}: Here $\Omega =
  [0, 
  \infty)$ and we assume that $f(t) = \beta_0 + \tau W_t$ where $W_t
  \sim BM$, $\beta_0 \sim N(0, C)$ and $\tau > 0$ (and independence
  between   $\beta_0$ and $\{W_t\}$). $C$ will be taken to be
  large. Now  
  \begin{equation*}
    K(s, t) = C + \tau^2 \min(s, t). 
  \end{equation*}
\item \textbf{Integrated Brownian Motion}: $\Omega = [0,
  \infty)$ and
  \begin{equation*}
    f(t) = \int_0^t W_s ds ~~ \text{where $W_s \sim BM$}. 
  \end{equation*}
  The kernel is given by
  \begin{align*}
    K(s, t) &= \text{Cov} \left(\int_0^s W_u du, \int_0^t W_v dv
              \right) \\
            &= \int_0^s \int_0^t \text{Cov}(W_u, W_v) dv du 
            = \int_0^s \int_0^t \min(u, v) dv du
  \end{align*}
  To simplify further, assume that $s \leq t$ so that
  \begin{align*}
K(s, t)            &= \int_0^s \left(\int_0^u v dv + \int_u^t u dv \right) du
    \\
            &= \int_0^s \left(\frac{u^2}{2} + u(t - u) \right) du = \int_0^s \left(ut - \frac{u^2}{2} \right) du = t \frac{s^2}{2}
      - \frac{s^3}{6}. 
  \end{align*}
  For general $s$ and $t$, we have
  \begin{equation*}
    K(s, t) = \frac{1}{2} \max(s, t) \left(\min(s, t)\right)^2 -
    \frac{1}{6} \left(\min(s, t) \right)^3. 
  \end{equation*}
\item \textbf{(scaled) Integrated Brownian Motion Plus a Linear Term}:
  In the 
  IBM model, we have $f(0) = 0$ and $f'(0) = 0$. This might be
  an unrealistic assumption to make when $f$ is completely
  unknown. In this case, a better model might be
  \begin{equation*}
    f(t) = \beta_0 + \beta_1 t + \tau \int_0^t W_s ds
  \end{equation*}
  where $\beta_0, \beta_1, \{W_s\}$ are independent with $W_s$ being
  Brownian motion and $\beta_0, \beta_1 \overset{\text{i.i.d}}{\sim} N(0,
  C)$. The kernel now becomes
  \begin{align*}
    K(s, t) &= \text{Cov} \left(\beta_0 + \beta_1 s + \tau \int_0^s W_u du,
      \beta_0 + \beta_1 t + \tau \int_0^t W_v dv \right) \\ &= C (1 +
                                                              st) +
                                                              \frac{\tau^2}{2}
                                                              \max(s,
                                                              t)
                                                              \left(\min(s,
                                                              t)\right)^2
                                                              - 
    \frac{\tau^2}{6} \left(\min(s, t) \right)^3. 
  \end{align*}
\end{enumerate}


%\begin{thebibliography}{99}

%\bibitem[MacKay and Peto(1995)]{MacKay1995HierarchicalDirichlet}
%D.~J.~C. MacKay and L.~C. Peto,
%\newblock ``A hierarchical Dirichlet language model,''
%\newblock \emph{Natural Language Engineering}, vol.~1,
%no.~3, pp.~289--308, 1995.

%\bibitem[Rubin(1981)]{Rubin1981BayesianBootstrap}
%D.~B. Rubin,
%\newblock ``The Bayesian bootstrap,''
%\newblock \emph{The Annals of Statistics}, vol.~9,
%no.~1, pp.~130--134, 1981.
  
%\bibitem[Fisher(1934)]{Fisher1934}
%R.~A. Fisher,
%\newblock ``Probability, likelihood and quantity of information in the logic of
%uncertain inference,''
%\newblock \emph{Proceedings of the Royal Society of London. Series~A,
%Containing Papers of a Mathematical and Physical Character}, vol.~146,
%no.~856, pp.~1--8, 1934.

%\bibitem[Efron(2010)]{Efron}
%Efron, B. (2010)
%\textit{Large-scale inference: empirical Bayes methods for estimation,
%testing, and prediction}. Cambridge University Press, 2010.

%\bibitem[Shumway and Stoffer(2017)]{ShumwayStoffer}
%Shumway, R. H. and Stoffer, D. S. (2017)
%\textit{Time Series Analysis and Its Applications: With R
%  Examples}. Springer, 4th edition. 
  
%     \bibitem[Jaynes(2003)]{Jaynes} Jaynes, E. T, (2003)
%  \textit{Probability theory: the logic of science}. Cambridge
%  University Press, 2003.
 
%  \bibitem[Sinai(1992)]{Sinai} Sinai, Y. G, (1992)
%  \textit{Probability theory: an introductory course}. Springer
%  Verlag, 1992.  

%\bibitem[{deGroot(1986)}]{DeGrootBlackwell}DeGroot, M. H. (1986)
%       A conversation with David Blackwell.
%\newblock \emph{Statistical Science}, \textbf{1}, 40--53.

%         \bibitem[Mosteller(1987)]{Mosteller} Mosteller, F, (1987)
%  \textit{Fifty challenging problems in probability with
%    solutions}. Courier Corporation, 1987.

%    \bibitem[MacKay(2003)]{MacKay} MacKay, D. J. C., (2003)
%  \textit{Information theory, inference, and learning
%  algorithms}. Cambridge University Press, 2003.

%\bibitem[{Lindley(2013)}]{Lindley}Lindley, D. V. (2013)
%   \textit{Understanding Uncertainty}. John Wiley and sons, 2013.


%\end{thebibliography}







\end{document}

%%% Local Variables:
%%% mode: latex
%%% TeX-master: t
%%% End:
